\chapter{Theoretical Background}
\label{cap:fundamentacao}

This chapter presents the concepts that support the remainder of this work.
Section~\ref{sec:fund-regressao} defines the regression task and the metrics
used to evaluate regression models. Section~\ref{sec:fund-algoritmos} describes
the five learning algorithms adopted in the experiments.
Section~\ref{sec:fund-pipeline} presents the pipeline steps investigated in this
study and the procedures used to estimate model performance.
Section~\ref{sec:fund-vazamento} discusses how data leakage may arise in each of
these steps. Finally, Section~\ref{sec:fund-estatistica} describes the
statistical tests used in the analysis of the results.

\section{Regression and Its Evaluation}
\label{sec:fund-regressao}

In supervised learning, a model is trained from a set of instances for which the
expected output is known. Each instance is described by a vector of input
features $\mathbf{x}$ and is associated with an output value $y$, called the
target. The learning algorithm uses these instances to build a model capable of
predicting the target of new instances. When the target is categorical, the task
is called classification; when it is a continuous numerical value, such as the
price of a property or the body fat percentage of a person, the task is called
regression \cite{hastie2009}. This work addresses regression tasks only.

The purpose of a model is to make accurate predictions for instances that were
not available during training. A model that fits the training data very closely
may still perform poorly on new data, a phenomenon known as overfitting.
Therefore, model performance must be measured on data that were not used for
training, and the result is taken as an estimate of the model's generalization
performance, that is, of how it is expected to behave in practice.

In regression, the quality of a prediction is measured by its deviation from the
true value. Given $m$ evaluation instances with true values $y_i$ and predicted
values $\hat{y}_i$, the four metrics adopted in this work are defined as:

\begin{align}
  \text{MAE}  &= \frac{1}{m}\sum_{i=1}^{m} \lvert y_i - \hat{y}_i \rvert, \label{eq:mae}\\
  \text{MSE}  &= \frac{1}{m}\sum_{i=1}^{m} (y_i - \hat{y}_i)^2, \label{eq:mse}\\
  \text{RMSE} &= \sqrt{\text{MSE}}, \label{eq:rmse}\\
  R^2         &= 1 - \frac{\sum_{i=1}^{m} (y_i - \hat{y}_i)^2}{\sum_{i=1}^{m} (y_i - \bar{y})^2}, \label{eq:r2}
\end{align}

\noindent where $\bar{y}$ is the mean of the true values. The Mean Absolute Error
(MAE) is the average magnitude of the errors. The Mean Squared Error (MSE)
squares each error before averaging, so that large errors have a greater
influence on the result. The Root Mean Squared Error (RMSE) is the square root of
the MSE and is expressed in the same unit as the target, which makes it easier
to interpret. For these three metrics, lower values indicate better models.
However, they depend on the scale of the target: the same RMSE may represent a
small error for a target measured in thousands and a large error for a target
measured in units. Consequently, their values cannot be directly compared across
data sets.

The coefficient of determination ($R^2$) does not have this limitation. It
measures the proportion of the variance of the target that is explained by the
model, using as reference a baseline model that always predicts the mean of the
target. A value of 1 corresponds to perfect predictions, a value of 0 indicates
that the model performs as well as the baseline, and negative values indicate
that it performs worse than the baseline. Because it is independent of the scale
of the target, $R^2$ allows comparisons across data sets and is considered more
informative than the other metrics for this purpose \cite{chicco2021}. For this
reason, $R^2$ is the primary metric of this work, and RMSE is used as a
complementary metric.

\section{Regression Algorithms}
\label{sec:fund-algoritmos}

Five algorithms were considered in this work. They were selected because they
represent different learning approaches, which may lead to different
sensitivities to data leakage.

\subsection{Linear Regression}

Linear Regression (LR) models the target as a linear combination of the
features, $\hat{y} = \beta_0 + \beta_1 x_1 + \dots + \beta_d x_d$. The
coefficients $\beta_j$ are estimated by Ordinary Least Squares (OLS), which
minimizes the sum of the squared errors on the training data \cite{hastie2009}.
LR is simple, computationally efficient, and has no hyperparameters in its basic
form. A property of LR that is relevant to this work is that its predictions do
not change when a feature is shifted or multiplied by a constant, since the
corresponding coefficient is adjusted accordingly. Some variants impose
constraints on the coefficients, such as requiring all of them to be
non-negative, and these options can be treated as hyperparameters.

\subsection{K-Nearest Neighbors}

The K-Nearest Neighbors (KNN) algorithm is an instance-based method: it stores
the training data and predicts the target of a new instance as the average
target of the $K$ training instances closest to it \cite{cover1967,hastie2009}.
Its main hyperparameters are the number of neighbors $K$ and the distance
function used to measure the proximity between instances, such as the
Euclidean or the Manhattan distance. Since distances are computed directly from
the feature values, features with larger ranges have a greater influence on the
choice of the neighbors than features with smaller ranges. For this reason, KNN
is highly sensitive to the scale of the features.

\subsection{Decision Trees}

A Decision Tree (DT) recursively divides the feature space into regions. At each
node, the algorithm selects a feature and a threshold that split the training
instances into two groups, choosing the split that most reduces the error of the
predictions in the resulting groups. This process is repeated until a stopping
criterion is met, and the prediction of each terminal node, or leaf, is the mean
target of the training instances that reach it \cite{breiman1984}. Since each
split only compares a feature with a threshold, decision trees are not affected
by the scale of the features. However, a tree that grows without restrictions
tends to overfit the training data, so its complexity is usually controlled by
hyperparameters such as the maximum depth.

\subsection{Random Forests}

A Random Forest (RF) is an ensemble of decision trees whose predictions are
averaged \cite{breiman2001}. Each tree is trained on a bootstrap sample of the
training instances and, at each node, considers only a random subset of the
features. This randomness makes the trees different from each other, and
averaging their predictions produces a model that is more stable than a single
tree. Like decision trees, random forests are not affected by the scale of the
features. Their main hyperparameters are the number of trees and the maximum
depth of each tree.

\subsection{Support Vector Regression}

Support Vector Regression (SVR) adapts the Support Vector Machine to regression
problems \cite{drucker1997,smola2004}. It searches for a function that deviates
from the true targets by at most $\varepsilon$ for as many training instances as
possible, while keeping the function as simple as possible. Errors smaller than
$\varepsilon$ are ignored, and larger errors are penalized according to the
regularization parameter $C$. Through the use of kernel functions, such as the
linear and the Radial Basis Function (RBF) kernels, SVR can also model non-linear
relations. Since both $\varepsilon$ and the kernel functions depend on the scale
of the data, SVR generally requires standardized features, and often a
standardized target as well, to perform properly.

\section{Building and Evaluating ML Pipelines}
\label{sec:fund-pipeline}

In practice, raw data are rarely provided directly to a learning algorithm.
They are usually transformed by a sequence of preparation steps, and the
algorithm must be configured and evaluated. This complete sequence, from the raw
data to the final performance estimate, is referred to as an ML pipeline. This
section describes the four pipeline steps investigated in this work and the
procedures used to evaluate a model.

\subsection{Feature Scaling}

The features of a data set frequently have very different ranges, for instance
when some of them are measured in units and others in thousands. Feature scaling
transforms the features so that they have comparable ranges. The most common
method, standardization (or $z$-score normalization), transforms each feature
$x_j$ as follows:
\begin{equation}
  x'_j = \frac{x_j - \mu_j}{\sigma_j},
  \label{eq:zscore}
\end{equation}
where $\mu_j$ and $\sigma_j$ are the mean and the standard deviation of the
feature. After the transformation, every feature has mean zero and unit standard
deviation. It is important to note that $\mu_j$ and $\sigma_j$ are estimated from
the data. The scaler, therefore, also learns parameters from data before it can
be applied, and the choice of which data are used for this purpose is relevant
for the discussion in Section~\ref{sec:fund-vazamento}.

\subsection{Missing Value Imputation}

Real-world data sets often contain missing values, caused, for example, by
measurement failures or by incomplete records. Since most learning algorithms
cannot handle missing values, they must be replaced by estimated values before
training, a process called imputation. The mechanisms that generate missing
values are commonly classified into three types \cite{rubin1976,little2019}:
missing completely at random (MCAR), when the probability of a value being
missing is unrelated to the data; missing at random (MAR), when it depends only
on other observed features; and missing not at random (MNAR), when it depends on
the missing value itself.

The simplest imputation methods replace each missing value with the mean or the
median of the observed values of the same feature. The KNN imputer, in turn,
finds the $K$ instances most similar to the incomplete one, based on the
observed features, and uses the average of their values \cite{troyanskaya2001}.
As in feature scaling, all these methods learn information from the data (means,
medians, or reference instances) before they can be applied.

\subsection{Feature Selection}

Feature selection consists of keeping only a subset of the most relevant
features and discarding the others. Reducing the number of features can produce
simpler and faster models and, in some cases, improve their predictive
performance \cite{guyon2003}. A common approach, adopted in this work, is to
assign a score to each feature according to its relation with the target and to
keep the features with the highest scores. In regression, a usual score is
obtained from the $F$-test, which measures the strength of the linear relation
between each feature and the target. Unlike scaling and imputation, this step
uses the values of the target to decide which features are kept.

\subsection{Hyperparameter Tuning}

Hyperparameters are parameters of an algorithm that are defined before training,
such as the number of neighbors in KNN, the maximum depth of a decision tree, or
the regularization parameter and the kernel in SVR. Hyperparameter tuning is the
process of evaluating different values and selecting those that produce the best
performance. The most common strategy is grid search, which evaluates every
combination of values in a predefined grid; random search, which evaluates a
random sample of combinations, is a frequent alternative \cite{bergstra2012}. In
both cases, each candidate configuration must be evaluated on data not used to
train it, so the tuning process also relies on one of the evaluation procedures
described next.

\subsection{Model Evaluation}

The simplest evaluation procedure is the hold-out method, in which the data are
split once into a training set, used to build the model, and a test set, used
only to measure its performance. The resulting estimate, however, depends on a
single random split, which makes it unreliable for small data sets. In $k$-fold
cross-validation (CV), the data are divided into $k$ subsets of similar size,
called folds. The model is trained $k$ times, each time using $k-1$ folds for
training and the remaining fold for testing, and the final estimate is the
average of the $k$ results \cite{kohavi1995,hastie2009}. In this way, every
instance is used exactly once for testing. The procedure can also be repeated
with different random partitions to obtain a more stable estimate.

When the same CV is used both to select the hyperparameters and to report the
performance of the selected model, the reported value is the best result among
many evaluated configurations. Part of this result is due to chance, and the
estimate is therefore optimistically biased \cite{varma2006,cawley2010}. Nested
CV addresses this problem by using two levels of CV. The outer level divides the
data into training and test folds, as in a standard CV. Within each outer
training fold, an inner CV is used to select the hyperparameters, and the
selected model is then evaluated on the corresponding outer test fold, which did
not take part in the selection. Although nested CV is computationally more
expensive, it provides a nearly unbiased estimate of the performance of the
complete tuning procedure \cite{varma2006}.

\section{Data Leakage in ML Pipelines}
\label{sec:fund-vazamento}

The evaluation procedures described above are based on a fundamental
assumption: the data used to evaluate a model must not be used, in any way, to
build it. Data leakage occurs when this assumption is violated, that is, when
information that would not be available when the model is applied in practice is
used during its development \cite{kaufman2012}. Some forms of leakage are easily
identified, such as a feature that directly encodes the target or the same
record appearing in both the training and the test sets. Other forms are more
subtle. \citeonline{kapoor2023} describe, among other types, the leakage caused
by the lack of a clean separation between training and test data during
pre-processing, for instance when a pre-processing step is applied to the whole
data set before it is split.

This form of leakage is the focus of this work. As discussed in
Section~\ref{sec:fund-pipeline}, scalers, imputers, feature selectors, and tuning
procedures all learn information from the data. In a correct pipeline, each of
these components is fitted using only the training data of each split and is
then applied, without modification, to the test data. In a leakage-prone
pipeline, the same component is fitted once on the complete data set, and only
afterwards are the data split. As a consequence, the test instances contribute
to building the model that they are supposed to evaluate.
Table~\ref{tab:fontes-vazamento} summarizes, for each step, what is learned from
the data and which information from the test set reaches the model when leakage
occurs.

\begin{table}[htb]
\centering
\caption{What each pipeline step learns from data and what leaks when it is fitted on the full data set.}
\label{tab:fontes-vazamento}
\small
\begin{tabular}{p{2.6cm}p{5.0cm}p{6.2cm}}
\toprule
\textbf{Step} & \textbf{Learned from data} & \textbf{Information leaked from the test set} \\
\midrule
Scaling & Mean and standard deviation of each feature & Distribution of the test feature values \\
Imputation & Mean, median, or nearest neighbors of each feature & Test feature values used to impute training values (and vice versa) \\
Feature selection & Relevance of each feature with respect to the target & Test \emph{target} values, through the choice of the features that are kept \\
Hyperparameter tuning & Best configuration in the search space & Test performance, through the choice of the configuration that is reported \\
\bottomrule
\end{tabular}
\\[2pt]
\footnotesize Source: The Author
\end{table}

The table suggests that the four steps are not equally harmful. Scaling and
imputation leak only information about the features. When training and test
instances come from the same distribution, this information is already well
estimated from the training data alone, and the test data add little to it.
Feature selection and hyperparameter tuning, on the other hand, leak information
related to the target, either directly or through the performance of the model.
This allows the pipeline to adapt to the specific test instances and may lead to
considerably overestimated performance. These expectations are evaluated
empirically in Chapter~\ref{cap:resultados}.

\section{Statistical Tests}
\label{sec:fund-estatistica}

In this work, each leakage-prone pipeline is compared with a leakage-free
pipeline executed on the same data split, which produces paired results. The
Wilcoxon signed-rank test \cite{wilcoxon1945} is a non-parametric test that
evaluates whether the differences between paired observations are
symmetrically distributed around zero. Since it is based on the ranks of the
differences rather than on their values, it does not assume that the differences
are normally distributed and is robust to extreme values. \citeonline{demsar2006}
recommends this test for comparing two methods over multiple data sets, using
one summary value per data set. Pairs with a difference of exactly zero are
discarded by the test, so a large number of ties reduces its effective sample
size.

To compare more than two methods over the same data sets, such as the five
algorithms considered in this work, the Friedman test \cite{friedman1937} can be
applied. It ranks the methods within each data set and tests whether their
average ranks differ more than would be expected by chance \cite{demsar2006}.

When many hypotheses are tested, the probability of obtaining at least one
significant result by chance increases. The Holm procedure \cite{holm1979}
adjusts the $p$-values to control this probability and is less conservative than
the classical Bonferroni correction. Finally, the Spearman correlation
coefficient $\rho$ \cite{spearman1904} measures the strength of a monotonic
relation between two variables. In this work, it is used to assess whether the
effect of leakage is related to characteristics of the data sets, such as their
number of instances.
